\documentclass[9pt,t,serif,professionalfont,xcolor={table,dvipsnames},noamsthm]{beamer}
\usepackage{xr-hyper}  % for cleveref; https://tex.stackexchange.com/a/244272/121234
\PassOptionsToPackage{pdfpagemode=FullScreen}{hyperref}
\setbeamersize{text margin left = 0.5cm,
               text margin right = 0.3cm}
\mode<presentation>
{
  \usetheme{boxes}
  \setbeamercovered{invisible}
  \setbeamertemplate{navigation symbols}{} 
}
\hypersetup{colorlinks=true,linkcolor=blue} 

% set the fonts
\usefonttheme{professionalfonts} % using non standard fonts for beamer
\usefonttheme{serif} % default family is serif
\usepackage{../presentation,../presentationfonts}

% load the pdfanim style, should be done after hyperref
% \usepackage{pdfanim}

% load and initialise animation
% arguments:
%   \PDFAnimLoad[options]{name}{xxx}{number}
%   - options: in this example the animation is looped
%   - name: name of the animation
%     (with this name it can be reused several times in the document)
%   - number: number of animation frames / files (n)
% the animation will be composed of files
% xxx0.pdf, xxx1.pdf ... xxx(n-1).pdf


\title{Definition of Computation}

\author{Jim Hef{}feron}
\institute{
  Mathematics and Statistics\\
  Saint Michael's College\\
  Colchester, VT\\[1ex]
  \texttt{hefferon.net}
}
\date{}

% This is only inserted into the PDF information catalog.
\subject{Introduction to Computability}

\usepackage{catchfilebetweentags} % read text from background.tex
\newcommand{\catchfilefn}{../../prologue/prologue.tex}
% Keep from swallowing end-of-lines
%   https://tex.stackexchange.com/a/40704/121234
\usepackage{etoolbox}
\makeatletter
\patchcmd{\CatchFBT@Fin@l}{\endlinechar\m@ne}{}
  {}{\typeout{Unsuccessful patch!}}
\makeatother

\usepackage{xr}
  \externaldocument{../../book}
  \externaldocument{../../prologue/prologue}

% Directory of this chapter's asy material
\newcommand{\asydir}{../../prologue/asy/}
\newcommand{\tmexercises}{\asydir tm_exercises/pdfs/}
% Directory of the src/asy material
\newcommand{\generalasydir}{../../asy/}

\usepackage{multirow,bigstrut}
\begin{document}
\begin{frame}
  \titlepage
\end{frame}




\section{Definition of computation}
% ------------------------------------------------------
\begin{frame}
  \frametitle{Introduction}

What can be done?

\pause
In 1928, D~Hibert and his student W~Ackermann proposed the
\textit{Entscheidungsproblem},
which asks for a definite procedure
(what we today call an algorithm) that decides, from an input mathematical
statement, whether that statement is true or false.

Being interested in what is true or false requires no
further explanation.
But why the definite procedure?

\pause
In Mathematics and Computer Science we take Euclid's development a paradigm.
The theorems follow from the axioms in a definite way: you can check,
mechanically, each step of the proofs.
So we start by defining a mechanical computation.
\end{frame}


\begin{frame}
  \frametitle{Turing machine}
A~Turing's approach to the \textit{Entscheidungsproblem} was to 
describe a mechanism that computes.
He imagined a clerk doing by-hand multiplication
with a sheet of paper that gradually becomes covered with columns of
numbers.
With that image he laid down a number of principles that a 
computing machine must satisfy.

\pause
\begin{itemize}
\item
\ExecuteMetaData[\catchfilefn]{discussion:DefTuringMachinei}
\pause\item
\ExecuteMetaData[\catchfilefn]{discussion:DefTuringMachineii}
\end{itemize}
\end{frame}

\begin{frame}
\begin{itemize}
\item
\ExecuteMetaData[\catchfilefn]{discussion:DefTuringMachineiii}
\pause\item
\ExecuteMetaData[\catchfilefn]{discussion:DefTuringMachineiv}
\end{itemize}
\end{frame}


\begin{frame}
Turing pictured a box containing a
mechanism and fitted with a tape.
\begin{center}
   \includegraphics{\generalasydir /share/machine00.pdf}
\end{center}
\ExecuteMetaData[\catchfilefn]{discussion:DefTuringMachinev}

\ExecuteMetaData[\catchfilefn]{discussion:DefTuringMachinevi}
\end{frame}

\begin{frame}
\ExecuteMetaData[\catchfilefn]{discussion:DefTuringMachinevii}

\ExecuteMetaData[\catchfilefn]{discussion:DefTuringMachineviii}

Before the definition, we first see a couple of examples.
\end{frame}


\begin{frame}{Computation of the predecessor}
\ExecuteMetaData[\catchfilefn]{ex:PredecessorTMi}

\ExecuteMetaData[\catchfilefn]{ex:PredecessorTMii}
\end{frame}

\begin{frame}
Here are the later steps of the computation.
\begin{equation*}
 \TM_{\text{pred}}=
  \set{ \tminstruction{0}{B}{L}{1},\,
        \tminstruction{0}{1}{R}{0},\,
        \tminstruction{1}{B}{L}{2},\,
        \tminstruction{1}{1}{B}{1},\,
        \tminstruction{2}{B}{R}{3},\,
        \tminstruction{2}{1}{L}{2} 
      }
\end{equation*}

\ExecuteMetaData[\catchfilefn]{ex:PredecessorTMiii}
\end{frame}

\begin{frame}
\ExecuteMetaData[\catchfilefn]{ex:PredecessorTMiv}
\begin{center}
  \begin{transitionfunction}
           &$\blank$      &\str{1}        \\ \cline{2-3}
     $q_0$ &$\str{L}q_1$  &$\str{R}q_0$     \\ 
     $q_1$ &$\str{L}q_2$  &$\str{B}q_1$     \\ 
     $q_2$ &$\str{R}q_3$  &$\str{L}q_2$     \\  
     $q_3$ &--            &--     \\
  \end{transitionfunction}
  \hspace{2em}
  \vcenteredhbox{\includegraphics{\asydir circlediagram/circlediagram00.pdf}}
\end{center}
\end{frame}



\begin{frame}{Computation of addition}
\ExecuteMetaData[\catchfilefn]{ex:AdditionTMi}
Here is the table and graph for $\TM_{\text{add}}$. 
\begin{center}
  \begin{transitionfunction}
           &$\blank$      &\str{1}        \\ \cline{2-3}
     $q_0$ &$\str{B}q_1$  &$\str{R}q_0$     \\ 
     $q_1$ &$\str{1}q_1$  &$\str{1}q_2$     \\ 
     $q_2$ &$\str{B}q_3$  &$\str{L}q_2$     \\  
     $q_3$ &$\str{R}q_3$  &$\str{B}q_4$     \\  
     $q_4$ &$\str{R}q_5$  &$\str{1}q_5$     \\  
     $q_5$ &--            &--     \\
  \end{transitionfunction}
  \hspace{2em}
  \vcenteredhbox{\includegraphics{\asydir circlediagram/circlediagram01.pdf}}
\end{center}
\ExecuteMetaData[\catchfilefn]{ex:AdditionTMii}
\end{frame}
\begin{frame}
Here are the the other steps in the computation
(the initialization is step~$0$).
\ExecuteMetaData[\catchfilefn]{ex:AdditionTMiii}
\end{frame}


\begin{frame}{Some Turing machines don't halt}
A crucial observation:~some Turing machines,
for at least some starting configurations, never halt.
\ExecuteMetaData[\catchfilefn]{ex:SomeTMsDontHalt}
% \begin{center}
%   \vcenteredhbox{\includegraphics{\asydir circlediagram/infloop.pdf}}
% \end{center}
\end{frame}


\begin{frame}{Definition of a Turing machine}
\ExecuteMetaData[\catchfilefn]{df:TM}

\medskip
\ExecuteMetaData[\catchfilefn]{df:TMnote}
\end{frame}


\begin{frame}{The action of a machine}
We can give a precise definition, a mathematical definition, 
of how Turing machines act.

\ExecuteMetaData[\catchfilefn]{def:TMConfigurationi}

\ExecuteMetaData[\catchfilefn]{def:TMConfigurationii}

\ExecuteMetaData[\catchfilefn]{def:TMConfigurationiii}

The remaining details are in the text.
The example below gives the idea.
\end{frame}
\begin{frame}
\begin{example}
This machine decides whether the input is an odd number.
\begin{equation*}
  \TM_{\text{odds}}=\set{
    \tminstruction{0}{1}{B}{1},\,
    \tminstruction{0}{B}{B}{4},\,
    \tminstruction{1}{B}{R}{2},\,
    \tminstruction{2}{1}{B}{3},\,
    \tminstruction{2}{B}{1}{4},\,
    \tminstruction{3}{B}{R}{0}
    }
\end{equation*}
We can trace through a computation one step at a time.
This machine starts with \str{111} on the tape, and the head pointing
to the leftmost~\str{1}.
\begin{align*}
  \sequence{q_0,\str{1},\emptystring,\str{11}}
  &\yields\sequence{q_1,\str{B},\emptystring,\str{11}}  \\
  &\yields\sequence{q_2,\str{1},\emptystring,\str{1}}  \\
  &\yields\sequence{q_2,\str{B},\emptystring,\str{1}}  \\
  &\yields\sequence{q_0,\str{1},\emptystring,\emptystring}  \\
  &\yields\sequence{q_1,\str{B},\emptystring,\emptystring}  \\
  &\yields\sequence{q_2,\str{B},\emptystring,\emptystring}  \\
  &\yields\sequence{q_4,\str{1},\emptystring,\emptystring}
\end{align*}
The machine now halts because there is no instruction for $q_pT_p=q_4\str{1}$.
\end{example}
\end{frame}



\begin{frame}{Definition of effective function}
\ExecuteMetaData[\catchfilefn]{df:FcnComputedByTMi}

\ExecuteMetaData[\catchfilefn]{df:FcnComputedByTMii}

\ExecuteMetaData[\catchfilefn]{df:FcnComputedByTMiii}

\ExecuteMetaData[\catchfilefn]{df:FcnConverge}
\end{frame}


\begin{frame}{Discussion: representations}
\ExecuteMetaData[\catchfilefn]{discussion:FcnConverge}

\begin{example}
We try not to worry too much about the representations.
If we have a Turing machine that, when given input that is a string, 
$\str{1}^n$, of $n$~many \str{1}'s, will output a string,
$\str{1}^{2n}$, of $2n$~many \str{1}'s, then we will feel free to
say that it computes the function $f(x)=2x$. 
\end{example}
\end{frame}


\begin{frame}
\ExecuteMetaData[\catchfilefn]{df:ComputableFcn}

\begin{example}
The computable functions form a very big collection.
All of these are computable.
\begin{equation*}
  f_0(x)=2x\qquad f_1(x)=x^2\qquad f_2(x)=2^x
\end{equation*}
So are these.
\begin{equation*}
  f_3(x)=
  \begin{cases}
    x/2 &\case{if $x$ is even}  \\
    3x+1 &\case{if $x$ is odd}
  \end{cases}
 \qquad
  f_4(x)=\text{the least prime number greater than x}
\end{equation*}
The computable sets are also a big collection.
These sets are computable.
\begin{equation*}
  S_0=\set{2n\suchthat n\in\N}
  \qquad
  S_1=\set{x\suchthat\text{the prime factorization has three distinct primes}}
\end{equation*}
\end{example}
\end{frame}


% ===========================================
\section{Church's Thesis}

\begin{frame}{Statement of Church's Thesis}
% \ExecuteMetaData[\catchfilefn]{ChurchsThesis}
\framebox{\alert{\parbox{0.95\textwidth}{The set of things that can be computed 
by a discrete and deterministic mechanism 
is the same as the set of things that can be computed by
a Turing machine.}}}
 
  Some authors call this the Church-Turing Thesis.  
  Here we figure that because Turing has the machine, 
  we can give Church sole possesion of the thesis.
\end{frame}

\begin{frame}{Evidence}
\begin{description}[xxx]  % optional argument sets width of indent
\item[Coverage]
Everything that people think of as intuitively computable has proven to
be computable on a Turing machine.
\item[Convergence]
There have been many models of computation proposed and all 
compute the same set of things, namely the set of things computed by
Turing machines.
\item[Consistency]
The details of the definition, including the number of tapes or tape
heads, the number of symbols available, the number of states, all
don't change the collection of things computed.
(This includes the condition of determinism.)
\item[Clarity]
Turing's analysis is persuasive of itself.
\end{description}
\end{frame}

\begin{frame}{An empirical question?}
As yet, there is no persuasive evidence of physically possible
systems that allow you to compute things that a Turing machine
cannot compute.

We shall take proposals involving wormholes, black hole event horizons,
etc., 
as too speculative to be as-yet persuasive.

In particular, note that the set of things computed by 
quantum computers is exactly the same as the set of things computed by 
Turing machines.
We expect that quantum computers will run faster, 
but are not more capable in principle.  
\end{frame}

\begin{frame}{Church's Thesis in arguments}
Church's Thesis equates `computable' with `computable on a Turing machine'.
We will leverage it in two ways.
\begin{enumerate}
\item It gives our results a larger, philosophical, importance.
When we produce a result \textit{No Turing machine can \ldots{}}, 
Church's Thesis changes it to \textit{This task cannot be done: \ldots{}}
\item To prove results about Turing machines we often 
argue intuitively.
For instance, consider the problem of whether this function is computable:~it 
inputs a number~$n$ and
returns~$1$ when $n$~is a prime, and~$0$ otherwise.
To show that can be computed by a Turing machine, we will not exhibit
such a machine.
We will instead argue that the function is intuitively computable, and invoke
Church's Thesis.

This is potentially dangerous. 
However, we have so much experience with computation
that our intuition is quite good.
For us, there is more of a danger
of being overwhelmed by detail than of being too vague. 
\end{enumerate}
\end{frame}



% ======================================================
\section{Primitive recursive functions}

\begin{frame}{Definition by recursion}
\ExecuteMetaData[\catchfilefn]{discussion:PrimReci}

\begin{example}
\ExecuteMetaData[\catchfilefn]{discussion:PrimRecii}
\end{example}
\end{frame}

\begin{frame}
\ExecuteMetaData[\catchfilefn]{discussion:PrimReciii}

\ExecuteMetaData[\catchfilefn]{discussion:PrimReciv}
\end{frame}

\begin{frame}{Defining by primitive recursion}
\ExecuteMetaData[\catchfilefn]{def:PrimRec}

\ExecuteMetaData[\catchfilefn]{ex:PrimReci}
\end{frame}

\begin{frame}
\ExecuteMetaData[\catchfilefn]{ex:PrimRecii}
\end{frame}

\begin{frame}{Collection of Primitive Recursive Functions}
\ExecuteMetaData[\catchfilefn]{def:PrimRecFcns}

The collection of primitive recursive functions is very large.
It includes addition,
proper subtraction, multiplication, integer division, exponentiation,
the equality Boolean function, the remainder function, the sum of a 
bounded number of terms, the product of a bounded number of terms, and
just about every function that comes to mind.
\end{frame}



% ===================================================
\section{General Recursive Functions}

\begin{frame}{Ackermann's function}
\ExecuteMetaData[\catchfilefn]{discussion:Ackermanni}

\ExecuteMetaData[\catchfilefn]{discussion:Ackermannii}
\end{frame}

\begin{frame}
\ExecuteMetaData[\catchfilefn]{discussion:Ackermanniii}
\end{frame}



\begin{frame}
\ExecuteMetaData[\catchfilefn]{def:Hyperoperation}
\ExecuteMetaData[\catchfilefn]{lem:Hyperoperation}
\ExecuteMetaData[\catchfilefn]{thm:Hyperoperation}

The hyperoperation is a form of \alert{Ackermann's function}.
\end{frame}


\begin{frame}{$\mu$ recursion}
\ExecuteMetaData[\catchfilefn]{def:MuRecursioni}
\ExecuteMetaData[\catchfilefn]{def:MuRecursionii}
% \end{frame}


% \begin{frame}
\begin{example}
A number~$n$ is a \alert{Carmichael number} if for all $b\in\set{1,2,\ldots,n}$
relatively prime to~$n$,
the number $b^{n-1}$ leaves a remainder of~$1$ on division by~$n$.
The first Carmichael number is~$561$.
The Boolean function
\begin{equation*}
  g(n)=
  \begin{cases}
    1  &\case{if $n$ is a Carmichael number}  \\
    0  &\case{otherwise}
  \end{cases}
\end{equation*}
is primitive recursive.
Here we define the Carmichael function $\map{C}{\N}{\N}$
by $\mu$-recursion.
\begin{equation*}
  C(n)=
  \begin{cases}
    561  &\case{if $n=0$}           \\
    \least y\,\big[ \text{$y>C(n-1)$ and $1-g(y)=0$} \big]
         &\case{otherwise}
    \end{cases}
\end{equation*}
\end{example}
\end{frame}



\begin{frame}
\ExecuteMetaData[\catchfilefn]{def:GeneralRecursive}

The collection of general recursive functions is the same set as the
collection of functions computed by some Turing machine, which we
call the `computable functions'.

The development of the computable functions as the recursive functions
has the advantage that we go straight to functions, without having to
worry about represenatations on a Turing machine.
It has the disadvantage that 
it is not evident that this is what can be mechanically computed
in the way that Turing's development is clear of itself.  
\end{frame}

\end{document}






